%% RESUMO EM LINGUA ESTRANGEIRA (elemento obrigatorio -- NBR 6028)
%% Traducao fiel do resumo em portugues. As keywords sao definidas em
%% config/dados.tex.

\begin{abstract}
This research investigates the application of spatiotemporal statistical models, machine learning techniques, and big data approaches to data from the National System of Construction Costs and Indices (SINAPI) in order to support decision-making in public infrastructure policies. The study is based on the recognition that SINAPI is the primary national reference for construction costs in Brazil and that its potential for predictive analytics remains underexplored. The research seeks to demonstrate the importance of forecasting SINAPI data and to assess whether the use of reference values from the São Paulo region for locations lacking local observations may lead to overpricing in public procurement contracts. The theoretical framework encompasses concepts from data science, artificial intelligence, machine learning, spatiotemporal analysis, and evidence-based public policy management. The proposed approach includes the development and validation of models capable of identifying trends, estimating future costs, and supporting budget planning for investments in housing, sanitation, and infrastructure. The study also examines the benefits of anticipating risks associated with budget insufficiency, project interruptions, and economic-financial imbalances in construction contracts. The data analyzed are derived from the public SINAPI database, maintained by CAIXA Econômica Federal in partnership with the Brazilian Institute of Geography and Statistics (IBGE), covering historical series of construction costs and inputs. As an expected outcome, this research aims to demonstrate that the application of advanced data analytics techniques can transform SINAPI from a predominantly consultative tool into an intelligence platform capable of supporting strategic decision-making, increasing the efficiency of public resource allocation, and strengthening evidence-based public policy formulation.
\end{abstract}
